\documentclass[
	ngerman,%globale Übergabe der Hauptsprache
%	logofile=example-image, %Falls die Logo Dateien nicht vorliegen
	authorontitle=true,
	]{bfhbeamer}


%\usepackage[main=ngerman]{babel}

% Der folgende Block ist nur bei pdfTeX auf Versionen vor April 2018 notwendig
\usepackage{iftex}
\ifPDFTeX
\usepackage[utf8]{inputenc}
\usepackage{blkarray}
\usepackage{graphicx}%kompatibilität mit TeX Versionen vor April 2018
\fi


%Makros für Formatierungen der Doku
%Im Allgemeinen nicht notwendig!
\let\code\texttt

\title{TraceSentry}
\subtitle{AI-Aided Caches-n-Logs Monitoring-n-Wiping Daemon}
\author[L. Ammann, J. Scherer, L. Scherer]{Luca Ammann, Janic Scherer, Luca Scherer}
\date{8. Januar 2025}
\institute{BFH-TI}
\titlegraphic*{\includegraphics{assets/presentation/title}}%is only used with BFH-graphic and BFH-fullgraphic

%Activate the output of a frame number:
%\setbeamertemplate{page number in head/foot}[framenumber]

\begin{document}

\setbeamertemplate{title page}[BFH-graphic]
\maketitle

\begin{frame}{Inhalt}
	\framesubtitle{Übersicht}
	\begin{columns}
		\begin{column}{0.4\textwidth}
			\tableofcontents
		\end{column}
		\begin{column}{0.6\textwidth}
			\begin{center}
				%\includegraphics[width=0.3\textwidth]
				%\includegraphics[width=0.6\textwidth]
			\end{center}
		\end{column}
	\end{columns}
\end{frame}


\section{Problemstellung}\label{sec:problemstellung}
\setbeamertemplate{Problemstellung}[BFH-ruled]
\frame{\sectionpage}

\begin{frame}{Problemstellung}
	\framesubtitle{Problem}
	\begin{itemize}
		\item Cache- und Logdateien: tempor\"are Dateien von Betriebssystemen und Anwendungen.
		\item Viele Benutzer wissen nicht:
		\begin{itemize}
			\item Welche Dateien existieren.
			\item Welche Anwendungen diese manipulieren.
		\end{itemize}
		\item Risiken: Datenschutzprobleme und fehlende Kontrolle.
	\end{itemize}
\end{frame}

\begin{frame}{Problemstellung}
	\framesubtitle{Chance}
	\begin{itemize}

		\item Entwickeln einer Software f\"ur:
		\begin{itemize}
			\item \textbf{Suche} von Dateien.
			\item \textbf{\"Uberwachung} von Dateien bzgl. deren Erstellung, Ver\"anderung und L\"oschung.
			\item \textbf{Analyse und Bewertung} von Dateien.
			\item \textbf{Bereinigung} von Dateien.
		\end{itemize}
		\item Nutzung eines KI-Modells f\"ur:
		\begin{itemize}
			\item Analyse von Struktur (Syntax) und Bedeutung (Semantik) von Dateien.
			\item Bewertung des Dateiinhalts bzgl.
			dessen Schädlichkeit.
			\item Empfehlungen von n\"otigen Massnahmen.
		\end{itemize}
	\end{itemize}
\end{frame}

\begin{frame}{Projektziel}
	\begin{itemize}
		\item Entwicklung einer \textbf{plattformunabh\"angigen, benutzerfreundlichen} Anwendung mit Dokumentation.
		\item Praxisnahe \textbf{Fallstudie} zur Evaluation des Potentials von aktuellen KI-Technologien.
		\item Fokus auf \textbf{Resultate} und Kernfunktionalit\"at.
	\end{itemize}
\end{frame}

\begin{frame}{Rahmenbedingungen}
	\begin{itemize}
		\item Entwicklung eines Hintergrundprozesses (\textbf{Daemon}).
		\item Kommandozeilen-Interface f\"ur Benutzerinteraktion (\textbf{CLI}).
		\item \textbf{Cache- und Log}dateien als prim\"arer Fokus, Erweiterung auf andere Dateitypen m\"oglich.
		\item \textbf{FLOSS-lizenzierte} und \textbf{plattformunabh\"angige} Software.
		\item Evaluation eines aktuellen KI-Modells. \textbf{Keine Eigenentwicklung}.
	\end{itemize}
\end{frame}



\section{Lösung}\label{sec:loesung}
\setbeamertemplate{Lösung}[BFH-ruled]
\frame{\sectionpage}

\begin{frame}{Der Daemon}
	\begin{columns}
		\begin{column}{0.7\textwidth}
			\begin{itemize}
				\item \textbf{Periodisch \& automatisch:} Ausführung von Überwachungsaufgaben
				\begin{itemize}
					\item Erstellen von Dateisystem-Snapshots (Hashbäume).
					\item Vergleich von Snapshots (mit Vor- bzw.\ Nachfolger).
				\end{itemize}
				\item \textbf{Via REST-API:} Kommunikation mit Benutzer
				\begin{itemize}
						  \item Suche von Dateien (z.B.\ mit regulären Ausdrücken).
						  \item Management von überwachten Pfaden.
						  \item Operationen auf Snapshots.
						  \item Anfragen von KI-Einschätzungen.
						  \item Löschen von Dateien.
						  \item Administrative Funktionen (Starten/Stoppen des Daemons).
				\end{itemize}
			\end{itemize}
		\end{column}
		\begin{column}{0.3\textwidth}
			\begin{center}
				\includegraphics[width=0.5\textwidth]{assets/presentation/ghost-smile}
			\end{center}
		\end{column}
	\end{columns}
	\note[item]{https://www.svgrepo.com/}

\end{frame}

\begin{frame}{Die CLI}
	\framesubtitle{Kommandozeilen-Schnittstelle}
	\begin{columns}
		\begin{column}{0.7\textwidth}
			\begin{itemize}
				\item \textbf ...
			\end{itemize}
		\end{column}
		\begin{column}{0.3\textwidth}
			\begin{center}
				\includegraphics[width=0.5\textwidth]{assets/presentation/prompt}
			\end{center}
		\end{column}
	\end{columns}
	\note[item]{...}

\end{frame}

\section{Demo}\label{sec:demo}
\setbeamertemplate{Demo}[BFH-ruled]
\frame{\sectionpage}

\begin{frame}{Demo Frame}
	\framesubtitle{Untertitel}
	Bla Bla
\end{frame}

\section{Projektmanagement}\label{sec:projektmanagement}
\setbeamertemplate{Projektmanagement}[BFH-ruled]
\frame{\sectionpage}

\begin{frame}{Projektmanagement Frame}
	\framesubtitle{Untertitel}
	Bla Bla
\end{frame}

\section{Fazit}\label{sec:fazit}
\setbeamertemplate{Fazit}[BFH-ruled]
\frame{\sectionpage}

\begin{frame}{Fazit Frame}
	\framesubtitle{Untertitel}
	Bla Bla
\end{frame}

\section{Ausblick}\label{sec:ausblick}
\setbeamertemplate{Ausblick}[BFH-ruled]
\frame{\sectionpage}

\begin{frame}{Ausblick Frame}
	\framesubtitle{Untertitel}
	Bla Bla
\end{frame}

\begin{frame}{Fragen Frame}
	\framesubtitle{Untertitel}
	Bla Bla
\end{frame}

\section{examples}\label{sec:examples}
\begin{frame}{Blocks}
	\begin{block}{Block with a title}
		Content.
	\end{block}
	\begin{block}{}
		Without title
	\end{block}
\end{frame}

\begin{frame}{Block types}
	\begin{exampleblock}{Exampleblock}
		Content.
	\end{exampleblock}
	\begin{alertblock}{Alertblock}
		Content.
	\end{alertblock}
	\begin{example}[Example environment]
		Content.
	\end{example}
\end{frame}

\end{document}

